\chapter{大数律、中心极限定理与渐近线性化}\label{chap:math-stat-asymptotics}

统计量的有限样本分布通常依赖整个总体分布，难以写成闭式；渐近理论则把问题分成两个不同尺度。首先，大数律说明经验平均 $P_ng$ 在 $O(1)$ 尺度上趋于确定极限 $Pg$；随后，中心极限定理描述其 $n^{-1/2}$ 量级的随机涨落。连续映射、Slutsky、Delta 方法与渐近线性表示只是把这两个基本极限传递到更复杂的统计量。后面的 $M$-估计、MLE、LAN、Wald、Score 与 BvM 都依赖这一层，因此本章不仅给结论，也明确哪些条件控制一致性、哪些条件只控制一阶正态近似。

\section{随机变量列的收敛与概率阶}

设 $X_n,X$ 定义在同一概率空间。若对任意 $\varepsilon>0$，
\[
 \mathbb P(|X_n-X|>\varepsilon)\to0,
\]
则记 $X_n\xrightarrow{\mathbb P} X$；若
\[
 \mathbb P\{\omega:X_n(\omega)\to X(\omega)\}=1,
\]
则记 $X_n\xrightarrow{a.s.} X$。若 $X_n$ 的分布弱收敛到 $X$ 的分布，则记 $X_n\xrightarrow{d} X$。在实线上，弱收敛等价于分布函数在 $F_X$ 的所有连续点收敛；在一般可分度量空间中，也可以用所有有界连续函数 $f$ 的积分
\[
 \mathbb E f(X_n)\to\mathbb E f(X)
\]
来刻画。

三种收敛满足
\[
 X_n\xrightarrow{a.s.} X
 \quad\Longrightarrow\quad
 X_n\xrightarrow{\mathbb P} X
 \quad\Longrightarrow\quad
 X_n\xrightarrow{d} X,
\]
逆向一般不成立。当极限是常数 $c$ 时，$X_n\xrightarrow{d} c$ 与 $X_n\xrightarrow{\mathbb P} c$ 等价，这一点在估计量一致性中经常使用。

为了记录随机误差的尺度，定义 $O_{\mathbb P}$ 与 $o_{\mathbb P}$。若 $X_n/a_n$ 在概率意义下有界，即对任意 $\varepsilon>0$ 存在 $M<\infty$ 使充分大 $n$ 时
\[
 \mathbb P(|X_n|>Ma_n)<\varepsilon,
\]
则记 $X_n=O_{\mathbb P}(a_n)$；若 $X_n/a_n\xrightarrow{\mathbb P}0$，则记 $X_n=o_{\mathbb P}(a_n)$。中心极限定理通常给出
\[
 \overline X_n-\mu=O_{\mathbb P}(n^{-1/2}),
\]
而一个二阶 Taylor 余项若是 $O_{\mathbb P}(\|\widehat\theta-\theta\|^2)$，在 $\sqrt n$ 尺度上往往就是 $o_{\mathbb P}(1)$。

若 $X_n\xrightarrow{d} X$ 且 $g$ 在 $X$ 几乎处处连续，则连续映射定理给出
\[
 g(X_n)\xrightarrow{d} g(X).
\]
若 $X_n\xrightarrow{d} X$、$Y_n\xrightarrow{\mathbb P} c$，则 Slutsky 定理给出
\[
 X_n+Y_n\xrightarrow{d} X+c,
 \qquad
 X_nY_n\xrightarrow{d} cX,
\]
并在 $c\ne0$ 时有
\[
 \frac{X_n}{Y_n}\xrightarrow{d}\frac{X}{c}.
\]
它的统计意义非常具体：一个未知尺度若能被一致估计，就可以在一阶渐近分布中用估计值替换；但这只说明渐近等价，并不自动产生有限样本精确分布。

\section{经验平均、大数律与统一收敛}

在证明大数律之前，先把“尾概率如何控制极限”写成统一工具。Markov 不等式从非负随机变量的一阶矩出发：
\[
P(Y\ge a)\le\frac{\mathbb E Y}{a}.
\]
取 $Y=(X-\mathbb E X)^2$ 得 Chebyshev；取 $Y=\exp(tX)$ 并优化 $t$ 则得到 Chernoff 方法。若独立变量 $X_i\in[a_i,b_i]$，Hoeffding 不等式给出
\begin{equation}
P\!\left(
\sum_{i=1}^n(X_i-\mathbb E X_i)\ge t
\right)
\le
\exp\!\left(
-\frac{2t^2}{\sum_i(b_i-a_i)^2}
\right).
\label{eq:math-stat-hoeffding}
\end{equation}
与 Chebyshev 的多项式尾界相比，它揭示 bounded independent sums 的指数集中。

若只假设 $|X_i-\mathbb E X_i|\le M$ 且总方差
\[
V_n=\sum_i\Var(X_i),
\]
Bernstein 不等式保留方差信息：
\begin{equation}
P\!\left(
\sum_i(X_i-\mathbb E X_i)\ge t
\right)
\le
\exp\!\left(
-\frac{t^2}{2(V_n+Mt/3)}
\right).
\label{eq:math-stat-bernstein}
\end{equation}
当 $t\ll V_n/M$ 时指数近似 Gaussian $\exp(-t^2/2V_n)$；当 $t$ 很大时转为 $\exp(-ct/M)$。这就是“中心区由方差控制、极端区由单项幅度控制”的两尺度结构。

\begin{derivation}[Chernoff 方法如何产生 Hoeffding 型指数界]
对任意 $\lambda>0$，Markov 不等式给出
\[
P(S_n\ge t)
=P(\ee^{\lambda S_n}\ge\ee^{\lambda t})
\le
\ee^{-\lambda t}\mathbb E\ee^{\lambda S_n}.
\]
若 $S_n=\sum_iY_i$ 且 $Y_i$ 独立中心化，则
\[
\mathbb E\ee^{\lambda S_n}
=\prod_i\mathbb E\ee^{\lambda Y_i}.
\]
Hoeffding lemma 对 $Y_i\in[a_i-\mathbb E X_i,b_i-\mathbb E X_i]$ 给出
\[
\log\mathbb E\ee^{\lambda Y_i}
\le\frac{\lambda^2(b_i-a_i)^2}{8}.
\]
其证明只用两点分布的凸性：把 $Y_i$ 视作区间 $[a_i,b_i]$ 上的随机变量，由 $\ee^{\lambda y}$ 关于 $y$ 的凸性，
\[
\ee^{\lambda Y_i}
\le
\frac{b_i-Y_i}{b_i-a_i}\ee^{\lambda a_i}
+\frac{Y_i-a_i}{b_i-a_i}\ee^{\lambda b_i},
\]
取期望并令 $p_i=(\mathbb E Y_i-a_i)/(b_i-a_i)$，得
\[
\mathbb E\ee^{\lambda Y_i}
\le
(1-p_i)\ee^{\lambda a_i}+p_i\ee^{\lambda b_i}
=\ee^{L(\lambda)},
\]
其中 $L(\lambda)=\lambda\mathbb E Y_i+\log\{(1-p_i)+p_i\ee^{\lambda(b_i-a_i)}\}$。中心化使 $\mathbb E Y_i=0$，再利用 $\log u\le u-1$ 与 $L''(\lambda)\le(b_i-a_i)^2/4$，对 $L$ 在 $0$ 处 Taylor 展开即得 Hoeffding lemma。
所以
\[
P(S_n\ge t)
\le
\exp\!\left[
-\lambda t
+\frac{\lambda^2}{8}\sum_i(b_i-a_i)^2
\right].
\]
右侧关于 $\lambda$ 最小值在
\[
\lambda_*
=\frac{4t}{\sum_i(b_i-a_i)^2}
\]
处取得，代回即得\eqnrefs{eq:math-stat-hoeffding}。因此 concentration inequality 的核心不是某个特殊分布，而是"指数矩 + 独立性 + 优化"三步。

若进一步保留方差信息 $v_i=\Var(Y_i)$，Bennett 不等式给出
\[
\log\mathbb E\ee^{\lambda Y_i}
\le
\frac{v_i}{M^2}\bigl(\ee^{\lambda M}-1-\lambda M\bigr),
\qquad |Y_i|\le M,
\]
对 $\lambda$ 优化并利用 $\ee^u-1-u\le u^2/(2(1-u/3))$ 即得\eqnrefs{eq:math-stat-bernstein}。这把 Hoeffding 中只依赖区间长度的界换成同时依赖方差与截断尺度的界，在稀疏信号或异方差情形下显著更紧。
\end{derivation}

举例：正态均值的有界截断。若 $X_i\iid\mathcal N(0,1)$ 但只观测 $Y_i=X_i\mathds{1}_{\{|X_i|\le M\}}$，则 $|Y_i|\le M$、$\Var(Y_i)=v\le1$。对 $S_n=\sum_i(Y_i-\mathbb E Y_i)$，Bernstein 给出
\[
\mathbb P(S_n\ge t)
\le
\exp\!\left\{-\frac{t^2}{2(nv+Mt/3)}\right\},
\]
当 $t\ll nv/M$ 时退化为 $\exp(-t^2/2nv)$，与 CLT 的 Gaussian 中心区一致；当 $t$ 进入大偏差区时转为 $\exp(-3t/2M)$，体现单项幅度控制。

与统计力学的联系。Chernoff 熵
\[
I(a)=\sup_{\lambda\ge0}\{\lambda a-\log\mathbb E\ee^{\lambda X}\}
\]
正是大偏差速率函数，与热力学 Legendre 变换 $\sup_\beta\{\beta E-\log Z(\beta)\}$ 同构：$\lambda$ 对应逆温度 $\beta$，$\log\mathbb E\ee^{\lambda X}$ 对应自由能 $\log Z$，$I(a)$ 对应比自由能。Cramér 定理 $P(\overline X_n\approx a)\asymp\ee^{-nI(a)}$ 因此是统计力学配分函数估计在概率论中的对应物。


Borel--Cantelli 引理把一列尾概率转成样本路径陈述。对事件 $A_n$，若
\[
\sum_nP(A_n)<\infty,
\]
则
\begin{equation*}P(A_n\ \text{i.o.})=0.\end{equation*}
这里不需要独立性。反之，若 $A_n$ 独立且 $\sum_nP(A_n)=\infty$，则
\begin{equation*}P(A_n\ \text{i.o.})=1.\end{equation*}
因此“依概率收敛”若能给出可求和的偏离概率，常可通过第一 Borel--Cantelli 直接升级成几乎处处收敛。强大数律的困难就在于朴素 Chebyshev 界通常只有 $O(1/n)$，其和发散，因此必须使用子序列、截断或更强集中工具。

设 $X_1,X_2,\ldots\iid P$。对可积函数 $g$，记
\[
 Pg=\mathbb E g(X_1),
 \qquad
 P_ng=\frac1n\sum_{i=1}^ng(X_i).
\]
于是大数律可以统一理解为经验测度 $P_n$ 对固定测试函数 $g$ 的作用趋于总体测度 $P$。

若 $\mathbb E|g(X_1)|<\infty$，iid 强大数律给出
\[
 P_ng\xrightarrow{a.s.} Pg,
\]
从而当然也有 $P_ng\xrightarrow{\mathbb P} Pg$。有限二阶矩不是结论成立的必要条件，但它给出最直接的弱大数律证明：若 $\Var(g(X_1))<\infty$，则
\[
 \Var(P_ng)=\frac{\Var(g(X_1))}{n},
\]
故 Chebyshev 不等式给出
\[
 \mathbb P(|P_ng-Pg|>\varepsilon)
 \le\frac{\Var(g(X_1))}{n\varepsilon^2}\to0.
\]

仅有有限一阶矩时，仍可通过截断控制重尾。这个结论之所以重要，是因为“样本平均稳定”依赖可积性，而不是依赖正态、矩母函数或有限方差。


为简化记号设 $g(X_i)=X_i$ 且 $\mathbb E|X_1|<\infty$。定义随 $n$ 改变的截断变量
\[
 Y_{n,i}=X_i\mathds{1}\{|X_i|\le n\},
 \qquad
 \overline Y_n=\frac1n\sum_{i=1}^nY_{n,i}.
\]
先控制截断事件。由尾部可积性，
\[
 \sum_{m=1}^\infty \mathbb P(|X_1|>m)
 \le \mathbb E|X_1|<\infty,
\]
因而 $m\mathbb P(|X_1|>m)\to0$。于是
\[
 \mathbb P(\overline X_n\ne\overline Y_n)
 \le n\mathbb P(|X_1|>n)\to0.
\]
同时
\[
 |\mathbb E Y_{n,1}-\mathbb E X_1|
 \le\mathbb E\{|X_1|\mathds{1}(|X_1|>n)\}\to0.
\]
对截断平均，
\[
 \Var(\overline Y_n)
 =\frac1n\Var(Y_{n,1})
 \le\frac1n\mathbb E\{X_1^2\mathds{1}(|X_1|\le n)\}.
\]
把最后一个期望按 $|X_1|\le n$ 分层，或用
\[
 \frac{X_1^2}{n}\mathds{1}(|X_1|\le n)
 \le |X_1|
\]
并结合逐点收敛到 $0$，由支配收敛定理得到
\[
 \Var(\overline Y_n)\to0.
\]
所以 Chebyshev 不等式给出
\[
 \overline Y_n-\mathbb E Y_{n,1}\xrightarrow{\mathbb P}0.
\]
将“原平均与截断平均之差”“截断平均的随机波动”“截断均值与原均值之差”三项相加，即得
\[
 \overline X_n\xrightarrow{\mathbb P}\mathbb E X_1.
\]

Cauchy 分布（密度 $1/(\pi(1+x^2))$）：$\mathbb E|X_1|=\infty$，弱大数律不适用——事实上 $\bar X_n$ 仍服从 Cauchy 分布，不收敛。Pareto 分布（密度 $\alpha x^{-\alpha-1}$，$x\ge1$）：$\alpha>1$ 时 $\mathbb E X_1<\infty$，$\bar X_n\xrightarrow{\mathbb P}\alpha/(\alpha-1)$；$\alpha\le1$ 时均值发散。截断技巧的关键在于：有限期望保证截断均值与原均值之差在截断水平趋于无穷时消失，而截断后随机变量方差有界，可用 Chebyshev 控制。


弱大数律只要求依概率控制；强大数律必须控制整条样本轨道。对 iid 情形，有限一阶矩本身已经足够：
\begin{equation}\label{eq:math-stat-slln-l1}
 \mathbb E|X_1|<\infty
 \quad\Longrightarrow\quad
 \overline X_n\xrightarrow{a.s.}\mathbb E X_1.
\end{equation}
证明不能简单把 Chebyshev 概率界逐项求和，因为在只有一阶矩时方差甚至可能不存在。正确结构是先截去极端观测，再把截断后的随机涨落变成一个可求和的方差级数。


记 $\mu=\mathbb E X_1$，并对每个指标按其自身尺度截断
\[
 Y_k=X_k\mathds{1}\{|X_k|\le k\}.
\]
由尾积分公式与尾概率的单调性，
\[
 \sum_{k=1}^\infty\mathbb P(X_k\ne Y_k)
 =\sum_{k=1}^\infty\mathbb P(|X_1|>k)
 \le\mathbb E|X_1|<\infty.
\]
Borel--Cantelli 引理因此给出
\[
 \mathbb P(X_k\ne Y_k\ \text{无穷多次})=0.
\]
所以几乎处处存在随机 $K$，使 $k\ge K$ 后恒有 $X_k=Y_k$；于是
\[
 \frac1n\sum_{k=1}^n(X_k-Y_k)\xrightarrow{a.s.}0.
\]

接下来控制截断变量的中心化涨落。由 $\Var(Y_k)\le\mathbb E(Y_k^2)$，
\begin{align*}
 \sum_{k=1}^\infty\frac{\Var(Y_k)}{k^2}
 &\le
 \sum_{k=1}^\infty
 \frac{\mathbb E\{X_1^2\mathds{1}(|X_1|\le k)\}}{k^2}\\
 &=\mathbb E\!\left[
 X_1^2
 \sum_{k\ge \lceil |X_1|\rceil\vee1}\frac1{k^2}
 \right].
\end{align*}
对 $y=|X_1|$，若 $0\le y<1$，括号内乘积至多为 $(\pi^2/6)y^2\le(\pi^2/6)y$；若 $y\ge1$，取 $m=\lceil y\rceil$，则
\[
 \sum_{k=m}^\infty\frac1{k^2}
 \le\frac1{m^2}+\int_m^\infty\frac{\dd x}{x^2}
 \le\frac2m
 \le\frac2y.
\]
因此存在绝对常数 $C$ 使
\[
 \sum_{k=1}^\infty\frac{\Var(Y_k)}{k^2}
 \le C\mathbb E|X_1|<\infty.
\]
令
\[
 Z_k=\frac{Y_k-\mathbb E Y_k}{k}.
\]
这些变量独立、均值为零，而且 $\sum_k\Var(Z_k)<\infty$。下面由 Kolmogorov 最大不等式推出级数 $\sum_kZ_k$ 几乎处处收敛。

对任意独立零均值变量 $Z_{m+1},\ldots,Z_n$ 与 $a>0$，令尾部分和
\[
 T_j^{(m)}=\sum_{k=m+1}^jZ_k.
\]
定义首次越过阈值的互斥事件
\[
 B_j
 =\left\{
 |T_\ell^{(m)}|<a\ \text{for }m<\ell<j,
 \ |T_j^{(m)}|\ge a
 \right\}.
\]
在 $B_j$ 上，$T_j^{(m)}$ 只由 $Z_{m+1},\ldots,Z_j$ 决定，而 $T_n^{(m)}-T_j^{(m)}$ 只由后续独立零均值增量决定，所以交叉项条件期望为零。于是
\begin{align*}
 \mathbb E\!\left[(T_n^{(m)})^2\mathds{1}_{B_j}\right]
 &=\mathbb E\!\left[(T_j^{(m)})^2\mathds{1}_{B_j}\right]
 +\mathbb E\!\left[(T_n^{(m)}-T_j^{(m)})^2\mathds{1}_{B_j}\right]\\
 &\ge a^2\mathbb P(B_j).
\end{align*}
对互斥的 $B_j$ 求和，并使用
\[
 \mathbb E(T_n^{(m)})^2=\sum_{k=m+1}^n\Var(Z_k),
\]
得到 Kolmogorov 最大不等式
\begin{equation}\label{eq:math-stat-kolmogorov-maximal}
 \mathbb P\!\left(
 \max_{m<j\le n}|T_j^{(m)}|\ge a
 \right)
 \le
 \frac1{a^2}\sum_{k=m+1}^n\Var(Z_k).
\end{equation}
因为方差级数收敛，可以选递增整数 $m_\nu$ 使
\[
 \sum_{k>m_\nu}\Var(Z_k)\le2^{-3\nu}.
\]
令 $n\to\infty$ 并在\eqnrefs{eq:math-stat-kolmogorov-maximal}中取 $a=2^{-\nu}$，得到
\[
 \mathbb P\!\left(
 \sup_{j>m_\nu}
 \left|\sum_{k=m_\nu+1}^jZ_k\right|>2^{-\nu}
 \right)
 \le2^{-\nu}.
\]
右端可求和，再次由 Borel--Cantelli，级数尾部几乎处处成为 Cauchy 列，因此
\[
 \sum_{k=1}^\infty\frac{Y_k-\mathbb E Y_k}{k}
\]
几乎处处收敛。

最后使用 Kronecker 引理。令
\[
 R_n=\sum_{k=1}^n\frac{a_k}{k}.
\]
若 $R_n$ 收敛，则分部求和给出
\[
 \frac1n\sum_{k=1}^na_k
 =R_n-\frac1n\sum_{k=1}^{n-1}R_k\longrightarrow0,
\]
因为 $R_n$ 与其 Ces\`aro 平均具有同一极限。令 $a_k=Y_k-\mathbb E Y_k$，得到
\[
 \frac1n\sum_{k=1}^n(Y_k-\mathbb E Y_k)\xrightarrow{a.s.}0.
\]
另一方面 $\mathbb E Y_k\to\mu$，故 Ces\`aro 平均给出
\[
 \frac1n\sum_{k=1}^n\mathbb E Y_k\to\mu.
\]
再加上最开始已经证明的 $n^{-1}\sum_{k\le n}(X_k-Y_k)\xrightarrow{a.s.}0$，最终得到
\[
 \overline X_n\xrightarrow{a.s.}\mu,
\]
即\eqnrefs{eq:math-stat-slln-l1}。整个证明中没有使用有限方差；二阶矩只出现在截断变量上，并通过一阶尾部可积性变成可求和的方差级数。



经验分布函数给出一个重要而特殊的统一大数律。它对应函数类
\[
 \mathcal F_{\rm cdf}
 =\{x\mapsto\mathds{1}\{x\le t\}:t\in\R\},
\]
这个函数类虽然由不可微的指示函数组成，却仍满足
\[
 \sup_{t\in\R}|(P_n-P)\mathds{1}\{x\le t\}|\xrightarrow{a.s.}0.
\]
这就是 \eqnrefs{eq:math-stat-glivenko-cantelli}。证明的关键不是参数连续性，而是分布函数的单调性把无限多个阈值压缩到有限网格。


先证明 $U_i\iid\mathrm{Unif}(0,1)$ 的情形，记经验分布为
\[
 G_n(u)=\frac1n\sum_{i=1}^n\mathds{1}\{U_i\le u\}.
\]
固定整数 $m\ge1$，取网格
\[
 u_j=\frac jm,
 \qquad j=0,1,\ldots,m.
\]
对每个固定 $j$，强大数律给出
\[
 G_n(u_j)\xrightarrow{a.s.} u_j.
\]
因为网格点只有有限个，所以对每个固定 $m$ 都存在概率为 $1$ 的事件 $\Omega_m$，在其上
\[
 \max_{0\le j\le m}|G_n(u_j)-u_j|\to0.
\]
真正要让网格宽度随后也趋于零，还必须把这些事件统一起来。令
\[
 \Omega_*=\bigcap_{m=1}^{\infty}\Omega_m.
\]
这是可数个概率一事件的交，因此 $\mathbb P(\Omega_*)=1$。下面固定任意 $\omega\in\Omega_*$。若 $u_j\le u\le u_{j+1}$，由 $G_n$ 单调不减，
\[
 G_n(u_j)\le G_n(u)\le G_n(u_{j+1}).
\]
于是
\begin{align*}
 G_n(u)-u
 &\le G_n(u_{j+1})-u_j
 \\&=\{G_n(u_{j+1})-u_{j+1}\}+\frac1m,
\end{align*}
同时
\begin{align*}
 u-G_n(u)
 &\le u_{j+1}-G_n(u_j)
 \\&=\{u_j-G_n(u_j)\}+\frac1m.
\end{align*}
因此
\[
 \sup_{0\le u\le1}|G_n(u)-u|
 \le
 \max_{0\le j\le m}|G_n(u_j)-u_j|+\frac1m.
\]
因此对每个整数 $m\ge1$，在同一个事件 $\Omega_*$ 上都有
\[
 \limsup_{n\to\infty}
 \sup_{0\le u\le1}|G_n(u)-u|
 \le\frac1m.
\]
再令 $m\to\infty$，上确界误差的极限上界只能为 $0$，从而
\[
 \sup_{0\le u\le1}|G_n(u)-u|\xrightarrow{a.s.}0.
\]

现在令一般分布函数为 $F$，取 $U_i\iid\mathrm{Unif}(0,1)$ 并令
\[
 X_i=F^{-1}(U_i),
 \qquad
 F^{-1}(u)=\inf\{x:F(x)\ge u\}.
\]
广义逆满足事件恒等式
\[
 \{F^{-1}(U_i)\le x\}
 =\{U_i\le F(x)\},
\]
故相应经验分布满足
\[
 F_n(x)=G_n\{F(x)\}.
\]
于是
\[
 \sup_{x\in\R}|F_n(x)-F(x)|
 \le
 \sup_{0\le u\le1}|G_n(u)-u|
 \xrightarrow{a.s.}0.
\]
这个论证不要求 $F$ 连续；原子分布的跳跃已经由广义逆自动处理。因而 Glivenko--Cantelli 是真正的分布无关一致收敛，而不是逐点 LLN 的形式改写。

独立平均只是遍历平均的一种最简单结构。对 martingale difference $D_k$ 满足
\[
\mathbb E(D_k\mid\mathcal F_{k-1})=0,
\]
若 $|D_k|\le c_k$，Azuma--Hoeffding 不等式给出
\begin{equation*}P\!\left(
\left|\sum_{k=1}^nD_k\right|\ge t
\right)
\le
2\exp\!\left(
-\frac{t^2}{2\sum_{k=1}^nc_k^2}
\right).\end{equation*}
所以很多具有适应性依赖的数据仍能保留 iid 和式相似的集中尺度。若更一般地
\[
\sum_{k=1}^\infty\frac{\mathbb E D_k^2}{k^2}<\infty,
\]
则 martingale convergence 与 Kronecker lemma 可推出
\[
\frac1n\sum_{k=1}^nD_k\to0
\quad\text{a.s.}
\]
这提供了一个比 iid 更宽的强大数律模板：关键不是“没有依赖”，而是新信息在给定过去后仍有零条件均值，并且累计二阶涨落受到控制。

\begin{derivation}[martingale difference 的强平均收敛]
设 $(D_k,\mathcal F_k)$ 是 martingale difference，即 $D_k\in L^2$ 为 $\mathcal F_k$-可测且
\[
\mathbb E(D_k\mid\mathcal F_{k-1})=0.
\]
令
\[
M_n=\sum_{k=1}^n\frac{D_k}{k}.
\]
先验证 $(M_n,\mathcal F_n)$ 是 martingale。由条件期望的线性，
\[
\mathbb E(M_n\mid\mathcal F_{n-1})
=M_{n-1}+\frac1n\mathbb E(D_n\mid\mathcal F_{n-1})
=M_{n-1}.
\]
再计算 $L^2$ 范数。对 $j<k$，用 tower property 与 $D_j$ 的 $\mathcal F_{k-1}$-可测性，
\[
\mathbb E(D_jD_k)
=\mathbb E\bigl[D_j\,\mathbb E(D_k\mid\mathcal F_{k-1})\bigr]
=0,
\]
故 $\{D_k/k\}$ 在 $L^2$ 中两两正交，从而
\[
\mathbb E M_n^2
=\sum_{k=1}^n\frac{\mathbb E D_k^2}{k^2}.
\]
由假设 $\sum_{k\ge1}\mathbb E D_k^2/k^2<\infty$，$(M_n)$ 为 $L^2$-bounded martingale。由 martingale 收敛定理，存在 $M_\infty\in L^2$ 使
\[
M_n\to M_\infty
\quad\text{a.s.}
\]
即级数 $\sum_{k\ge1}D_k/k$ 几乎处处收敛。应用 Kronecker lemma：若 $b_k\uparrow\infty$ 且 $\sum_{k\ge1}a_k/b_k$ 收敛，则
\[
\frac1{b_n}\sum_{k=1}^n b_k a_k\to0.
\]
取 $a_k=D_k$，$b_k=k$，得
\[
\frac1n\sum_{k=1}^n D_k\to0
\quad\text{a.s.}
\]
这把“方差可求和”转成了路径上的平均收敛。
\end{derivation}


参数估计往往不是对一个固定函数取平均，而是在许多参数值之间比较经验准则。例如
\[
 M_n(\theta)=P_nm_\theta,
 \qquad
 M(\theta)=Pm_\theta.
\]
若只知道对每个固定 $\theta$ 有 $M_n(\theta)\xrightarrow{\mathbb P} M(\theta)$，并不足以保证经验最大点逼近总体最大点，因为最大化操作会在整个参数空间中搜索。因此需要统一大数律
\begin{equation}\label{eq:math-stat-ulln}
 \sup_{\theta\in\Theta}|M_n(\theta)-M(\theta)|\xrightarrow{\mathbb P}0.
\end{equation}
一个常用充分条件是：$\Theta$ 紧，$m_\theta(x)$ 关于 $\theta$ 连续，并存在可积包络 $G(x)$ 使
\[
 \sup_{\theta\in\Theta}|m_\theta(x)|\le G(x),
 \qquad
 PG<\infty.
\]
更一般的统一收敛理论由 Glivenko--Cantelli 函数类描述。


补上可测性条件，并假设对几乎每个 $x$，$\theta\mapsto m_\theta(x)$ 连续。对 $\delta>0$ 定义局部振荡包络
\[
 w_\delta(x)
 =\sup_{\|\theta-\theta'\|<\delta}
 |m_\theta(x)-m_{\theta'}(x)|.
\]
固定 $x$ 后，紧集上的连续函数一致连续，所以
\[
 w_\delta(x)\to0
 \qquad(\delta\downarrow0).
\]
又有 $w_\delta\le2G$，且 $PG<\infty$，故支配收敛定理给出
\[
 Pw_\delta\to0.
\]
给定 $\varepsilon>0$，先取 $\delta$ 使 $Pw_\delta<\varepsilon$。由 $\Theta$ 的紧性，可取有限 $\delta$-网
\[
 \theta_1,\ldots,\theta_N,
\]
使任意 $\theta$ 都满足某个 $\|\theta-\theta_j\|<\delta$。于是
\begin{align*}
 \sup_{\theta\in\Theta}
 |(P_n-P)m_\theta|
 &\le
 \max_{1\le j\le N}|(P_n-P)m_{\theta_j}|
 +P_nw_\delta+Pw_\delta.
\end{align*}
有限个固定函数分别由 LLN 得到
\[
 \max_j|(P_n-P)m_{\theta_j}|\xrightarrow{\mathbb P}0,
\]
同时 $P_nw_\delta\xrightarrow{\mathbb P} Pw_\delta$。因此
\[
 \limsup_{n\to\infty}
 \sup_{\theta}|(P_n-P)m_\theta|
 \le 2Pw_\delta
 <2\varepsilon
\]
以概率趋于 $1$ 成立。再令 $\varepsilon\downarrow0$，得到\eqnrefs{eq:math-stat-ulln}。这个证明清楚区分了两层机制：有限网格上的随机误差由普通 LLN 控制，网格之间的连续参数误差由可积振荡包络控制。

经验分布函数 $F_n(t)=\frac1n\sum\mathbf 1_{X_i\le t}$：取 $\mathcal F=\{\mathbf 1_{(-\infty,t]}:t\in\mathbb R\}$，ULLN 给出 $\sup_t|F_n(t)-F(t)|\xrightarrow{\mathbb P}0$——这就是 Glivenko--Cantelli 定理。核密度估计 $\hat f_n(x)=\frac1{nh_n}\sum K((x-X_i)/h_n)$：对函数类 $\{K_h(x-\cdot):x\in\mathbb R\}$ 施加 ULLN 给出逐点一致收敛，是非参数统计的基础。


如果 $M(\theta)$ 在真值 $\theta_0$ 唯一最大，并且离开任意邻域后都存在正的目标函数间隙，即对每个 $\varepsilon>0$，
\[
 M(\theta_0)
 -\sup_{\|\theta-\theta_0\|\ge\varepsilon}M(\theta)>0,
\]
则任何满足
\[
 M_n(\widehat\theta_n)
 \ge \sup_{\theta\in\Theta}M_n(\theta)-o_{\mathbb P}(1)
\]
的近似最大化序列都有 $\widehat\theta_n\xrightarrow{\mathbb P}\theta_0$。这就是后面 $M$-估计和 MLE 一致性的基本 argmax 逻辑：先保证经验目标对总体目标的一致逼近，再利用总体目标的可识别极值。


总体目标 $M$ 在 $\theta_0$ 唯一最大，离开邻域后有正间隙 $\Delta_\varepsilon$。统一收敛保证经验目标 $M_n$ 逼近 $M$ 到 $\Delta_\varepsilon/4$ 以内，此时 $\theta_0$ 处的经验值比外部区域高至少 $\Delta_\varepsilon/2$，经验极大点不可能落在外部。

定义间隙。固定 $\varepsilon>0$，记
\begin{equation*}
 \Delta_\varepsilon
 =M(\theta_0)
 -\sup_{\|\theta-\theta_0\|\ge\varepsilon}M(\theta)>0.
\end{equation*}
这是总体目标在 $\theta_0$ 与外部区域之间的正间隙，由可识别性假设保证。

统一收敛事件。在事件
\begin{equation*}
 \sup_\theta|M_n(\theta)-M(\theta)|<\frac{\Delta_\varepsilon}{4}
\end{equation*}
上，经验目标一致逼近总体目标。由\eqnrefs{eq:math-stat-ulln}，此事件概率趋于 $1$。

外部区域上界。若 $\|\theta-\theta_0\|\ge\varepsilon$，则
\begin{equation*}
 M_n(\theta)
 \le M(\theta)+\frac{\Delta_\varepsilon}{4}
 \le M(\theta_0)-\Delta_\varepsilon+\frac{\Delta_\varepsilon}{4}
 =M(\theta_0)-\frac{3\Delta_\varepsilon}{4}.
\end{equation*}

$\theta_0$ 处下界。同时
\begin{equation*}
 M_n(\theta_0)
 \ge M(\theta_0)-\frac{\Delta_\varepsilon}{4}.
\end{equation*}

比较。两式相减，
\begin{equation*}
 M_n(\theta_0)-M_n(\theta)
 \ge\frac{3\Delta_\varepsilon}{4}+\frac{\Delta_\varepsilon}{4}
 =\frac{\Delta_\varepsilon}{2}>0,
 \qquad \|\theta-\theta_0\|\ge\varepsilon.
\end{equation*}
即 $\theta_0$ 处的经验值比外部区域高至少 $\Delta_\varepsilon/2$。若近似最大化误差为 $o_{\mathbb P}(1)$，充分大 $n$ 时该误差小于 $\Delta_\varepsilon/4$，经验极大点不可能落在外部区域。

结论。在统一收敛事件（概率趋于 $1$）与近似最大化误差充分小（概率趋于 $1$）的交上，$\|\widehat\theta_n-\theta_0\|<\varepsilon$。因此
\begin{equation*}
 \mathbb P(\|\widehat\theta_n-\theta_0\|\ge\varepsilon)\to0.
\end{equation*}
由 $\varepsilon>0$ 任意，$\widehat\theta_n\xrightarrow{\mathbb P}\theta_0$。


经验分布的一致收敛只描述 $F_n-F$ 的 $O(1)$ 极限；若要研究置信带、Kolmogorov--Smirnov 统计量或由整个经验分布构造的泛函，就必须进一步研究它的 $n^{-1/2}$ 随机涨落。定义经验过程
\begin{equation*}\alpha_n(x)
:=\sqrt n\bigl(F_n(x)-F(x)\bigr).\end{equation*}
对任意有限组点 $x_1,\ldots,x_m$，向量
\[
\bigl(
\alpha_n(x_1),\ldots,
\alpha_n(x_m)
\bigr)
\]
是 iid 向量平均的 $\sqrt n$ 缩放，因此由多元 CLT 收敛到均值为零的 Gaussian vector。协方差可直接算出
\[
\begin{aligned}
\Cov\bigl(
\mathds{1}_{(-\infty,x]}(X),
\mathds{1}_{(-\infty,y]}(X)
\bigr)
&=F(x\wedge y)-F(x)F(y).
\end{aligned}
\]
所以极限过程应具有 covariance kernel
\begin{equation}
\boxed{
K_F(x,y)
=F(x\wedge y)-F(x)F(y)
}.
\label{eq:math-stat-brownian-bridge-covariance}
\end{equation}
这正是标准 Brownian bridge $B(t)$ 经 $t=F(x)$ 复合后的协方差。

Donsker 定理把有限维 CLT 升级成函数空间弱收敛：在适当的有界函数空间中，
\begin{equation}
\boxed{
\alpha_n
\Longrightarrow
B\circ F
}.
\label{eq:math-stat-donsker-cdf}
\end{equation}
其中 $B$ 是 $[0,1]$ 上标准 Brownian bridge，满足
\[
B(0)=B(1)=0,
\qquad
\Cov(B(s),B(t))=s\wedge t-st.
\]
这里真正困难的部分不是有限维正态性，而是 tightness：必须排除随机函数把质量逃到越来越细的尺度。也正因为如此，functional CLT 不能由“对每个 $x$ 都应用普通 CLT”机械得到。

若 $F$ 连续，则 probability integral transform 使
\[
U_i:=F(X_i)\sim\mathrm{Unif}(0,1),
\]
经验过程可以还原到统一分布模型。于是
\[
\sup_x\sqrt n|F_n(x)-F(x)|
\Longrightarrow
\sup_{0\le t\le1}|B(t)|.
\]
这就是 Kolmogorov--Smirnov 统计量具有 distribution-free 极限的原因：极限中原本的未知 $F$ 被单调时间变换吸收到 $t=F(x)$，只剩同一个 Brownian bridge。

经验过程语言还能统一一大类“对所有函数同时”的涨落。给定函数类 $\mathcal G$，定义
\[
\mathbb G_n g
:=\sqrt n(P_ng-Pg),
\qquad g\in\mathcal G.
\]
若 $\mathbb G_n$ 在 $\ell^\infty(\mathcal G)$ 中弱收敛到 tight Gaussian process $\mathbb G_P$，就称 $\mathcal G$ 为 $P$-Donsker class。其 covariance 为
\[
\Cov(\mathbb G_Pf,\mathbb G_Pg)
=P(fg)-Pf\,Pg.
\]
单个函数的 CLT 只对应 $\mathcal G=\{g\}$；有限维多元 CLT 对应有限函数类；真正的 empirical process theory 研究无限函数类在 entropy、covering number 或 VC 维条件下何时仍能获得统一 Gaussian 极限。

\begin{derivation}[Brownian bridge covariance 从 empirical CDF 直接得到]
记
\[
Y_i(x):=\mathds{1}_{\{X_i\le x\}}.
\]
则
\[
F_n(x)=\frac1n\sum_{i=1}^nY_i(x),
\qquad
\mathbb E Y_i(x)=F(x).
\]
因此
\[
\alpha_n(x)
=\sqrt n\{F_n(x)-F(x)\}
=\frac1{\sqrt n}\sum_{i=1}^n
\bigl(Y_i(x)-F(x)\bigr).
\]
对两个点 $x,y$，独立性使不同 $i$ 的交叉协方差消失：当 $i\neq j$ 时 $Y_i(x)$ 与 $Y_j(y)$ 独立，故
\[
\Cov\bigl(Y_i(x)-F(x),Y_j(y)-F(y)\bigr)=0,
\]
于是
\[
\begin{aligned}
\Cov(\alpha_n(x),\alpha_n(y))
&=\frac1n\sum_{i=1}^n\Cov\bigl(Y_i(x),Y_i(y)\bigr)\\
&=\Cov(Y_1(x),Y_1(y))\\
&=\mathbb E\bigl[Y_1(x)Y_1(y)\bigr]-F(x)F(y).
\end{aligned}
\]
而指示函数乘积满足
\[
Y_1(x)Y_1(y)
=\mathds{1}_{\{X_1\le x\}}\mathds{1}_{\{X_1\le y\}}
=\mathds{1}_{\{X_1\le x\wedge y\}},
\]
其中 $x\wedge y=\min(x,y)$，这是因为两个不等式 $X_1\le x$、$X_1\le y$ 同时成立等价于 $X_1$ 不超过较小者。取期望得
\[
\mathbb E[Y_1(x)Y_1(y)]
=F(x\wedge y),
\]
得到\eqnrefs{eq:math-stat-brownian-bridge-covariance}。这个 covariance 对每个 $n$ 都已经成立；Donsker 定理真正新增的是"整个随机函数的联合极限"而非 covariance 计算本身。
\end{derivation}

Brownian bridge 的零端点结构。由 $F(-\infty)=0$、$F(+\infty)=1$，协方差 $\sigma(x,y)=F(x\wedge y)-F(x)F(y)$ 在两端满足
\[
\sigma(-\infty,y)=0,
\qquad
\sigma(+\infty,y)=F(y)-F(y)=0,
\]
故极限过程 $\mathbb B\circ F$ 在 $t=0,1$ 处a.s. 为零，这正是"bridge"命名的来源：标准 Brownian 运动 $\mathbb B(t)$ 在两端被条件 $\mathbb B(0)=\mathbb B(1)=0$ 拉回成桥。

举例：均匀分布的经验过程。若 $F(t)=t$ 在 $[0,1]$ 上，则
\[
\sigma(s,t)=s\wedge t-st,
\qquad 0\le s,t\le1,
\]
即标准 Brownian bridge 的协方差。对分位数泛函 $T(F)=F^{-1}(p)$，其 Hadamard 导数为
\[
\dot T_F(h)=-\frac{h(F^{-1}(p))}{f(F^{-1}(p))}.
\]
因此
\[
\sqrt n\{F_n^{-1}(p)-F^{-1}(p)\}
\xrightarrow{d}
-\frac{\mathbb B(p)}{f(F^{-1}(p))},
\]
方差为 $p(1-p)/f^2(F^{-1}(p))$，即 Bahadur 量化的渐近方差。可见 Brownian bridge 协方差同时控制经验 CDF 与经验分位数的渐近涨落。


functional delta method 进一步把\eqnrefs{eq:math-stat-donsker-cdf}传递给依赖整个分布的统计泛函。若 $T$ 在 $F$ 处具有 Hadamard 导数 $\dot T_F$，则形式上
\[
\sqrt n\{T(F_n)-T(F)\}
\Longrightarrow
\dot T_F(B\circ F).
\]
普通有限维 Delta 方法因此只是这一函数空间线性化的特例。分位数、累积分布函数泛函、某些 $M$-估计量和半参数 influence function 都可以放在这一统一框架中。

经验过程给出函数层面的联合涨落；把函数类缩到有限维，便回到中心极限定理与多元随机涨落。

大数律只说明误差趋于零，没有告诉它以什么速度缩小。若 $X_i\iid P$，
\[
 \mathbb E X_1=\mu,
 \qquad
 0<\sigma^2=\Var(X_1)<\infty,
\]
则 Lindeberg--L\'evy 中心极限定理给出
\begin{equation}\label{eq:math-stat-iid-clt}
 \frac{\sqrt n(\overline X_n-\mu)}{\sigma}
 \xrightarrow{d}\mathcal N(0,1).
\end{equation}
因此样本均值误差的自然尺度是 $n^{-1/2}$。这不是说 $\overline X_n$ “变成正态变量”，而是说把趋于零的误差放大 $\sqrt n$ 后出现非退化正态极限。

图~\figref{fig:clt-hist} 数值演示了中心极限定理：即便总体分布远非正态（均匀或指数），$n=12$ 个独立变量标准化求和后，直方图也已逼近标准正态密度。

\begin{figure}[htbp]
\centering
\includegraphics[width=0.92\textwidth]{clt_hist.pdf}
\caption{中心极限定理数值演示。左：$\mathrm{Uniform}(0,1)$ 标准化求和；右：$\mathrm{Exp}(1)$ 标准化求和；$n=12$，均已近似 $N(0,1)$（红线）。}
\label{fig:clt-hist}
\end{figure}



令
\[
 Y_i=\frac{X_i-\mu}{\sigma},
 \qquad
 Z_n^{\rm CLT}=\frac1{\sqrt n}\sum_{i=1}^nY_i.
\]
这个量就是样本均值误差在中心极限定理尺度上的标准化涨落。因为 $\mathbb E Y_i=0$、$\mathbb E Y_i^2=1$，特征函数在原点满足
\[
 \varphi_Y(t)
 =\mathbb E\ee^{itY_i}
 =1+it\,\mathbb E Y_i-\frac{t^2}{2}\mathbb E Y_i^2+o(t^2)
 =1-\frac{t^2}{2}+o(t^2).
\]
这里 $o(t^2)$ 项的存在性由 $\mathbb E Y_i^2<\infty$ 保证；若进一步有 $\mathbb E|Y_i|^{2+\delta}<\infty$，则余项可写成 $O(|t|^{2+\delta})$，对应 Berry--Esseen 速率 $n^{-\delta/2}$。
独立性给出
\[
 \varphi_{Z_n^{\rm CLT}}(t)
 =\prod_{i=1}^n\varphi_{Y_i/\sqrt n}(t)
 =\left\{\varphi_Y\!\left(\frac{t}{\sqrt n}\right)\right\}^n
 =\left\{1-\frac{t^2}{2n}+o(n^{-1})\right\}^n.
\]
取对数，并利用 $\log(1+u)=u+O(u^2)$（当 $|u|\to0$），
\[
 \log\varphi_{Z_n^{\rm CLT}}(t)
 =n\log\left\{1-\frac{t^2}{2n}+o(n^{-1})\right\}
 =n\left\{-\frac{t^2}{2n}+o(n^{-1})\right\}
 \to-\frac{t^2}{2}.
\]
所以
\[
 \varphi_{Z_n^{\rm CLT}}(t)\to\ee^{-t^2/2},
\]
右端是标准正态特征函数。L\'evy 连续性定理遂给出\eqnrefs{eq:math-stat-iid-clt}。

为什么是 $\sqrt n$ 而不是别的尺度。若改用 $n^\alpha$ 归一化，则
\[
\varphi_{n^{-\alpha}\sum Y_i}(t)
=\left\{1-\frac{t^2}{2n^{2\alpha}}+o(n^{-2\alpha})\right\}^n.
\]
当 $\alpha>1/2$ 时指数 $n\cdot n^{-2\alpha}\to0$，极限为 $1$（大数律尺度，涨落消失）；当 $\alpha<1/2$ 时指数发散，特征函数在 $t\ne0$ 处趋于 $0$，分布逃向无穷；只有 $\alpha=1/2$ 给出非退化正态极限。$\sqrt n$ 尺度因此由二阶矩唯一确定。

举例：Rademacher 之和。若 $Y_i=\pm1$ 等概率，则 $\varphi_Y(t)=\cos t$，于是
\[
\varphi_{Z_n^{\rm CLT}}(t)=\cos^n\!\left(\frac{t}{\sqrt n}\right)
=\left\{1-\frac{t^2}{2n}+O(n^{-2})\right\}^n
\to\ee^{-t^2/2},
\]
这是离散到连续的最简例子。Lindeberg 条件
\[
\frac1{n}\sum_{i=1}^n\mathbb E\bigl[Y_i^2\mathds{1}_{\{|Y_i|>\varepsilon\sqrt n\}}\bigr]
\to0
\]
在此自动成立，因为当 $n>1/\varepsilon^2$ 时指示函数恒为零；这正是"中心区由二阶矩控制、尾部贡献可忽略"的 Lindeberg 哲学。


独立同分布不是 CLT 的逻辑必要条件。对三角阵 $\{X_{n,k}:1\le k\le k_n\}$，设每行独立且中心化，记
\[
 s_n^2=\sum_{k=1}^{k_n}\Var(X_{n,k}).
\]
若对每个 $\varepsilon>0$，Lindeberg 条件
\[
 \frac1{s_n^2}
 \sum_{k=1}^{k_n}
 \mathbb E\!
 \left[X_{n,k}^2\mathds{1}\{|X_{n,k}|>\varepsilon s_n\}\right]
 \to0
\]
成立，则
\[
 \frac1{s_n}\sum_{k=1}^{k_n}X_{n,k}
 \xrightarrow{d}\mathcal N(0,1).
\]
Lindeberg 条件的意义是：归一化总方差不能由少数极端项主导。若存在某个 $\delta>0$ 使 Lyapunov 条件
\[
 \frac{1}{s_n^{2+\delta}}
 \sum_{k=1}^{k_n}\mathbb E|X_{n,k}|^{2+\delta}\to0,
\]
则它足以推出 Lindeberg 条件。事实上，在事件 $|X_{n,k}|>\varepsilon s_n$ 上有
\[
 |X_{n,k}|^2
 \le
 \frac{|X_{n,k}|^{2+\delta}}
 {(\varepsilon s_n)^\delta},
\]
所以
\begin{align*}
 &\frac1{s_n^2}
 \sum_{k=1}^{k_n}
 \mathbb E\!\left[
 X_{n,k}^2\mathds{1}\{|X_{n,k}|>\varepsilon s_n\}
 \right]\\
 &\qquad\le
 \frac1{\varepsilon^\delta s_n^{2+\delta}}
 \sum_{k=1}^{k_n}\mathbb E|X_{n,k}|^{2+\delta}
 \longrightarrow0.
\end{align*}
因此 Lyapunov 条件只是一个更容易核验、但更强的尾部充分条件。后面的异方差线性模型和独立非同分布得分都可以在这类结构下建立渐近正态性。

Lindeberg 条件还蕴含 Feller infinitesimality：
\[
\max_{k\le k_n}\frac{\Var(X_{n,k})}{s_n^2}\to0.
\]
它排除了任何单个变量占据总方差的有限比例。反过来，在 Feller 条件下，Lindeberg 条件对三角阵 CLT 接近精确条件；Lindeberg--Feller 定理因此比“有限二阶矩 + iid”更接近中心极限定理真正的结构核心。


设三角阵满足 Lindeberg 条件且每个 $X_{n,k}$ 已中心化，于是 $\Var(X_{n,k})=\mathbb E X_{n,k}^2$。固定 $\varepsilon>0$，对每个 $k$ 用示性函数把全空间划分成 $|X_{n,k}|\le\varepsilon s_n$ 与 $|X_{n,k}|>\varepsilon s_n$ 两段，方差随之分解为
\begin{align*}
\mathbb E X_{n,k}^2
&=\mathbb E\{X_{n,k}^2\mathds{1}(|X_{n,k}|\le\varepsilon s_n)\}\\
&\quad+\mathbb E\{X_{n,k}^2\mathds{1}(|X_{n,k}|>\varepsilon s_n)\}.
\end{align*}
在第一段的示性集合上处处有 $|X_{n,k}|\le\varepsilon s_n$，从而 $X_{n,k}^2\le\varepsilon^2s_n^2$，取期望不放大上界，故
\begin{align*}
\mathbb E\{X_{n,k}^2\mathds{1}(|X_{n,k}|\le\varepsilon s_n)\}
\le\varepsilon^2s_n^2\,\mathbb E\mathds{1}(|X_{n,k}|\le\varepsilon s_n)
\le\varepsilon^2s_n^2.
\end{align*}
第二段恰是 Lindeberg 尾项在指标 $k$ 处的取值，它不超过把所有 $j$ 的尾项相加后的总和，即
\begin{align*}
\mathbb E\{X_{n,k}^2\mathds{1}(|X_{n,k}|>\varepsilon s_n)\}
\le\sum_{j=1}^{k_n}
\mathbb E\{X_{n,j}^2\mathds{1}(|X_{n,j}|>\varepsilon s_n)\}.
\end{align*}
把两项估计代回方差分解，两端除以 $s_n^2$，得到对每个 $k$ 一致成立的界
\begin{align*}
\frac{\Var(X_{n,k})}{s_n^2}
\le\varepsilon^2
+\frac1{s_n^2}\sum_{j=1}^{k_n}
\mathbb E\{X_{n,j}^2\mathds{1}(|X_{n,j}|>\varepsilon s_n)\}.
\end{align*}
上式右端与 $k$ 无关，因此对左端取 $\max_{k\le k_n}$ 仍保持同一上界：
\begin{align*}
\max_{k\le k_n}\frac{\Var(X_{n,k})}{s_n^2}
\le\varepsilon^2
+\frac1{s_n^2}\sum_{j=1}^{k_n}
\mathbb E\{X_{n,j}^2\mathds{1}(|X_{n,j}|>\varepsilon s_n)\}.
\end{align*}
由 Lindeberg 条件，固定 $\varepsilon>0$ 时右端第二项当 $n\to\infty$ 趋于零，故
\begin{align*}
\limsup_{n\to\infty}\max_{k\le k_n}\frac{\Var(X_{n,k})}{s_n^2}
\le\varepsilon^2.
\end{align*}
上式对任意 $\varepsilon>0$ 同时成立，令 $\varepsilon\downarrow0$ 即得
\begin{align*}
\lim_{n\to\infty}\max_{k\le k_n}\frac{\Var(X_{n,k})}{s_n^2}=0,
\end{align*}
即 Feller infinitesimality：没有任何单项方差能占据总方差的有限比例。


若 $X_i\in\R^m$ iid，$\mathbb E X_i=\mu$，协方差矩阵 $\Sigma$ 有限，则多元 CLT 为
\begin{equation*}\sqrt n(\overline X_n-\mu)
 \xrightarrow{d}\mathcal N_m(0,\Sigma).
\end{equation*}
其核心不是额外计算多元密度，而是 Cram\'er--Wold 装置：对任意固定 $a\in\R^m$，
\[
 a^{\mathsf T}\sqrt n(\overline X_n-\mu)
 \xrightarrow{d}\mathcal N(0,a^{\mathsf T}\Sigma a),
\]
而所有线性投影共同确定联合极限分布。

若三阶绝对中心矩 $m_3=\mathbb E|X_1-\mu|^3<\infty$，Berry--Esseen 界进一步给出
\[
 \sup_x\left|
 \mathbb P\!\left\{\frac{\sqrt n(\overline X_n-\mu)}{\sigma}\le x\right\}
 -\Phi(x)
 \right|
 \le C\frac{m_3}{\sigma^3\sqrt n}.
\]
这区分了两个问题：CLT 只回答极限是否正态；有限样本近似有多准确，还依赖偏斜、尾部和所关注的概率区域。

中心极限定理刻画线性平均；一般统计量则要先用 Delta 方法或渐近线性表示化为这种平均。

设
\[
 \sqrt n(T_n-\theta)
 \xrightarrow{d}\mathcal N_d(0,\Sigma),
\]
且 $g:\R^d\to\R^q$ 在 $\theta$ 可微，Jacobian 为 $Dg(\theta)$。Taylor 展开给出
\[
 g(T_n)-g(\theta)
 =Dg(\theta)(T_n-\theta)+r_n,
\]
其中 $\|r_n\|=o_{\mathbb P}(\|T_n-\theta\|)$。因此
\begin{equation*}\sqrt n\{g(T_n)-g(\theta)\}
 \xrightarrow{d}
 \mathcal N_q\!\left(0,
 Dg(\theta)\Sigma Dg(\theta)^{\mathsf T}\right).
\end{equation*}
若一阶导数在真值处退化，必须保留下一阶。例如一维 $g'(\theta)=0$、$g''(\theta)\ne0$ 时，若 $\sqrt n(T_n-\theta)\xrightarrow{d} Z$，则
\[
 n\{g(T_n)-g(\theta)\}
 \xrightarrow{d}\frac12g''(\theta)Z^2.
\]
这说明“先做局部线性化，再对高斯极限取二次函数”会自然产生卡方型极限。Wald 统计量使用的正是后一机制：约束函数先由一阶 Delta 方法得到渐近正态，再把标准化后的高斯向量取平方长度；它并不要求约束函数的一阶导数退化。

更有组织力的表述是\emph{渐近线性表示}。若某个 $d$ 维估计量满足
\begin{equation}\label{eq:math-stat-asymptotically-linear}
 \sqrt n(\widehat\theta_n-\theta_0)
 =\frac1{\sqrt n}\sum_{i=1}^n\mathrm{IF}_{\theta_0}(X_i)
 +o_{\mathbb P}(1),
 \qquad
 \mathbb E_{\theta_0}\mathrm{IF}_{\theta_0}(X)=0,
\end{equation}
则 $\mathrm{IF}_{\theta_0}$ 称为一阶影响函数。若其二阶矩有限，则 iid 向量 $\mathrm{IF}_{\theta_0}(X_i)$ 的均值为零、协方差为 $V_{\rm IF}$，多元 CLT 给出
\[
 \frac1{\sqrt n}\sum_{i=1}^n\mathrm{IF}_{\theta_0}(X_i)
 \xrightarrow{d}\mathcal N_d(0,V_{\rm IF}).
\]
再与\eqnrefs{eq:math-stat-asymptotically-linear}中的 $o_{\mathbb P}(1)$ 余项应用 Slutsky 定理，得到
\[
 \sqrt n(\widehat\theta_n-\theta_0)
 \xrightarrow{d}\mathcal N_d(0,V_{\rm IF}),
 \qquad
 V_{\rm IF}=\mathbb E_{\theta_0}[\mathrm{IF}_{\theta_0}\mathrm{IF}_{\theta_0}^{\mathsf T}].
\]
后面估计方程、MLE 与稳健 sandwich 协方差都可以归约到\eqnrefs{eq:math-stat-asymptotically-linear}。

若 $\widehat V_n\xrightarrow{\mathbb P} V_{\rm IF}$，则 Slutsky 定理允许把未知协方差用一致估计替换。例如标量情形令 $V_{\rm IF}=\sigma_{\rm IF}^2>0$，并取 $\widehat\sigma_n^2\xrightarrow{\mathbb P}\sigma_{\rm IF}^2$，则
\[
 \frac{\sqrt n(\widehat\theta_n-\theta_0)}{\widehat\sigma_n}
 \xrightarrow{d}\mathcal N(0,1).
\]
这就是 studentization 的一般渐近逻辑。它与正态样本中的精确 $t$ 分布必须区分：一般总体下分母替换只产生渐近标准正态；只有球对称高斯投影提供有限样本下的独立性，才会得到精确 Student $t$。

样本分位数给出一个重要的非光滑例子：它的渐近正态性并不需要先假定 Bahadur 表示。设
\[
 \xi_p=F^{-1}(p),\qquad 0<p<1,
\]
并假设 $F$ 在 $\xi_p$ 的某个开邻域内连续可微，且 $f(\xi_p)>0$。这个条件比最弱条件略强，但它保证分位点局部唯一、邻域内无原子，并允许直接使用局部逆映射。令经验分位数
\[
 \widehat\xi_p=F_n^{-1}(p).
\]
则
\begin{equation}\label{eq:math-stat-quantile-asymptotic-normal}
 \sqrt n(\widehat\xi_p-\xi_p)
 \xrightarrow{d}
 \mathcal N\!\left(0,
 \frac{p(1-p)}{f(\xi_p)^2}\right).
\end{equation}
密度出现在分母有直接几何意义：若 $F$ 在分位点附近很平，同样大小的概率误差经逆映射会被放大成更大的位置误差。


固定 $t\in\R$，令
\[
 y_n=\xi_p+\frac{t}{\sqrt n},
 \qquad
 q_n=F(y_n).
\]
由 $F$ 在 $\xi_p$ 处可微，
\[
 q_n
 =p+\frac{f(\xi_p)t}{\sqrt n}
 +o(n^{-1/2}).
\]
广义逆的定义给出事件的精确等价关系
\[
 \{\widehat\xi_p\le y_n\}
 =\{F_n(y_n)\ge p\}.
\]
写
\[
 N_n(y_n)=\sum_{i=1}^n\mathds{1}(X_i\le y_n),
\]
则对每个 $n$，
\[
 N_n(y_n)\sim\mathrm{Bin}(n,q_n).
\]
若 $k_n=\lceil np\rceil$，便有
\begin{align*}
 \mathbb P\!\left\{
 \sqrt n(\widehat\xi_p-\xi_p)\le t
 \right\}
 &=\mathbb P\{N_n(y_n)\ge k_n\}\\
 &=\mathbb P\!\left\{
 \frac{N_n(y_n)-nq_n}{\sqrt{nq_n(1-q_n)}}
 \ge
 \frac{k_n-nq_n}{\sqrt{nq_n(1-q_n)}}
 \right\}.
\end{align*}
由于 $q_n\to p\in(0,1)$，Bernoulli 三角阵满足 Lindeberg 条件，因此
\[
 \frac{N_n(y_n)-nq_n}{\sqrt{nq_n(1-q_n)}}
 \xrightarrow{d}\mathcal N(0,1).
\]
另一方面，$0\le k_n-np<1$，故
\begin{align*}
 \frac{k_n-nq_n}{\sqrt{nq_n(1-q_n)}}
 &=\frac{k_n-np}{\sqrt{nq_n(1-q_n)}}
 -\frac{\sqrt n(q_n-p)}{\sqrt{q_n(1-q_n)}}\\
 &\longrightarrow
 -\frac{f(\xi_p)t}{\sqrt{p(1-p)}}.
\end{align*}
于是
\[
 \mathbb P\!\left\{
 \sqrt n(\widehat\xi_p-\xi_p)\le t
 \right\}
 \longrightarrow
 \Phi\!\left(
 \frac{f(\xi_p)t}{\sqrt{p(1-p)}}
 \right),
\]
这正是\eqnrefs{eq:math-stat-quantile-asymptotic-normal}。关键操作是把“位置小于某点”反演成一个 Bernoulli 计数事件。

渐近正态性还可以加强为渐近线性表示，但这一步需要比单点 CLT 更强的\emph{局部经验过程控制}。对每个固定 $M<\infty$，若
\begin{equation}\label{eq:math-stat-local-empirical-equicontinuity}
 \sup_{|t|\le M}
 \sqrt n\left|
 \{F_n-F\}\!\left(\xi_p+\frac{t}{\sqrt n}\right)
 -\{F_n-F\}(\xi_p)
 \right|
 \xrightarrow{\mathbb P}0,
\end{equation}
则由\eqnrefs{eq:math-stat-quantile-asymptotic-normal}先知
$\sqrt n(\widehat\xi_p-\xi_p)=O_{\mathbb P}(1)$，再在这个随机 $n^{-1/2}$ 邻域内使用\eqnrefs{eq:math-stat-local-empirical-equicontinuity}，得到更强的 Bahadur 表示
\begin{equation}\label{eq:math-stat-quantile-bahadur}
 \sqrt n(\widehat\xi_p-\xi_p)
 =-\frac{\sqrt n\{F_n(\xi_p)-p\}}{f(\xi_p)}
 +o_{\mathbb P}(1).
\end{equation}
对半直线指示函数这一 VC 类，\eqnrefs{eq:math-stat-local-empirical-equicontinuity}确实成立；但它是经验过程的局部增量定理，而不是普通一维 Taylor 公式的自动后果。把这项输入单独列出，可以清楚区分两个层次：分位数渐近正态性只需局部计数 CLT，而 Bahadur 线性化还要求随机分布函数在收缩邻域内足够平滑。

\begin{derivation}[局部经验过程输入如何产生 Bahadur 表示]
由连续性与 $f(\xi_p)>0$，样本分位数以概率趋于 $1$ 落在无原子的局部邻域中。因而
\[
 F_n(\widehat\xi_p^- )\le p\le F_n(\widehat\xi_p),
 \qquad
 F_n(\widehat\xi_p)-p=O(n^{-1}).
\]
把左式第二个误差记为 $r_n=O_{\mathbb P}(n^{-1})$，则
\begin{align*}
 r_n
 &=F_n(\widehat\xi_p)-p\\
 &=\{F_n(\xi_p)-p\}
 +\{F(\widehat\xi_p)-F(\xi_p)\}\\
 &\quad+
 \Bigl[
 \{F_n-F\}(\widehat\xi_p)
 -\{F_n-F\}(\xi_p)
 \Bigr].
\end{align*}
由\eqnrefs{eq:math-stat-quantile-asymptotic-normal}，
\[
 \widehat\xi_p-\xi_p=O_{\mathbb P}(n^{-1/2}).
\]
因此\eqnrefs{eq:math-stat-local-empirical-equicontinuity}把最后一项控制为
\[
 o_{\mathbb P}(n^{-1/2}).
\]
同时局部可微性给出
\[
 F(\widehat\xi_p)-F(\xi_p)
 =f(\xi_p)(\widehat\xi_p-\xi_p)
 +o_{\mathbb P}(n^{-1/2}).
\]
代回并注意 $r_n=o_{\mathbb P}(n^{-1/2})$，得到
\[
 f(\xi_p)(\widehat\xi_p-\xi_p)
 =-\{F_n(\xi_p)-p\}
 +o_{\mathbb P}(n^{-1/2}).
\]
乘以 $\sqrt n/f(\xi_p)$ 即得\eqnrefs{eq:math-stat-quantile-bahadur}。由此影响函数可读为
\[
 \mathrm{IF}_{\xi_p}(x)
 =\frac{p-\mathds{1}(x\le\xi_p)}{f(\xi_p)},
\]
其方差恰为 $p(1-p)/f(\xi_p)^2$，与\eqnrefs{eq:math-stat-quantile-asymptotic-normal}一致。
\end{derivation}

本章建立了后续统计理论真正反复使用的四个接口：统一 LLN 控制经验准则，一维与多元 CLT 控制随机一阶项，Delta 方法传播光滑变换，渐近线性表示把复杂估计量还原为 iid 平均。下一章暂时离开 $n\to\infty$，研究正态模型中为什么可以在每一个有限 $n$ 下由正交投影得到更强的精确抽样分布。

\section{CLT 的数值验证与应用实例}\label{sec:clt-examples}
Berry--Esseen 定理给出 CLT 的收敛速率上界：
\[
\sup_x\left|F_n(x)-\Phi(x)\right|\le\frac{C\,\mathbb E|X-\mu|^3}{\sigma^3\sqrt n}
\]
其中 $C\le0.4748$（Shevtsova 2011）。这给出有限样本下正态近似的最大偏差。

\begin{exbox}[均匀分布的 CLT 收敛]
$X\sim\mathrm{Uniform}(0,1)$，$\mu=1/2$，$\sigma^2=1/12$，$\mathbb E|X-\mu|^3=1/32$。Berry--Esseen 界：
\[
\sup_x|F_n(x)-\Phi(x)|\le\frac{0.4748\times1/32}{(1/\sqrt{12})^3\sqrt n}\approx\frac{0.612}{\sqrt n}
\]
对 $n=10$：上界 0.194，实际最大偏差约 0.052；$n=100$：上界 0.061，实际约 0.016。Berry--Esseen 界保守约 3--4 倍，但给出了保证收敛的显式速率。
\end{exbox}
\begin{exbox}[方差估计的 Delta 方法]
样本方差 $S^2$ 的渐近分布：$\sqrt{n}(S^2-\sigma^2)\xrightarrow{d} N(0,\mu_4-\sigma^4)$，其中 $\mu_4=\mathbb E(X-\mu)^4$。对正态 $X$，$\mu_4=3\sigma^4$，给出 $\sqrt{n}(S^2-\sigma^2)\xrightarrow{d} N(0,2\sigma^4)$。

对标准差 $S=\sqrt{S^2}$，Delta 方法（$g(x)=\sqrt{x}$，$g'(\sigma^2)=1/(2\sigma)$）：
\[
\sqrt{n}(S-\sigma)\xrightarrow{d} N\!\left(0,\frac{\mu_4-\sigma^4}{4\sigma^2}\right)
\]
正态情形：$N(0,\sigma^2/2)$。实例：$\sigma=2$，$n=50$ 时 $S$ 的标准误 $\approx\sigma/\sqrt{2n}=2/\sqrt{100}=0.2$。
\end{exbox}

\begin{exbox}[比值估计的 Delta 方法]
估计比值 $\theta=\mu_X/\mu_Y$（如平均收入/平均家庭规模）。Delta 方法给出
\[
\sqrt{n}(\hat\theta-\theta)\xrightarrow{d} N(0,\sigma_\theta^2),\qquad
\sigma_\theta^2=\frac{\sigma_X^2}{\mu_Y^2}-\frac{2\mu_X\sigma_{XY}}{\mu_Y^3}+\frac{\mu_X^2\sigma_Y^2}{\mu_Y^4}
\]
$\mu_X=50000$，$\mu_Y=2.5$，$\sigma_X=30000$，$\sigma_Y=1.0$，$\sigma_{XY}=5000$，$n=1000$：$\sigma_\theta\approx420$，95\% 置信区间半宽 $1.96\times420/\sqrt{1000}\approx26$，即 $\hat\theta=20000\pm26$。
\end{exbox}
\begin{exbox}[中位数的渐近分布]
对 $N(\mu,\sigma^2)$ 样本，中位数 $\tilde X$ 的渐近分布：
\[
\sqrt{n}(\tilde X-\mu)\xrightarrow{d} N\!\left(0,\frac{\sigma^2\pi}{2}\right)
\]
因为 $f(\mu)=1/(\sigma\sqrt{2\pi})$，方差 $1/(4f(\mu)^2)=\pi\sigma^2/2$。对 $\sigma=1$，$n=100$：中位数标准误 $\sqrt{\pi/200}\approx0.125$，而样本均值标准误 $1/\sqrt{100}=0.1$。均值效率更高（正态分布下），但中位数对离群值稳健。
\end{exbox}
